\section{Cahier des charges et fonctionnement général du cycle de développement}
%-------------------------------------------------------------------------------
\subsection{Contenu du stage}
\label{sub:Contenu du stage}
Le sujet principal de mon stage concernait ainsi la refonte de l'Espace Recruteur du site cadremploi.fr.
L'application actuelle se présente comme un espace réservé aux professionnels afin de les aider dans leur campagnes de recrutement.
Grâce à ce service de e-commerce, il leur est possible d'accéder à tous les produits et services qu'offre Cadremploi et de payer directement en ligne avec leur carte de crédit ou par chèque après réception de la facture.
Cette application web avait déjà fait l'objet d'un développement il y a plusieurs années et une refonte totale en était nécessaire, de manière à proposer des services nouveaux, un design plus actuel et des éléments de maintenabilité plus importants.
Concrètement, les principaux objectifs de ce stage étaient les suivants:
\begin{itemize}
  \item{} Développer de nouvelles fonctionnalités dans un environnement technique dynamique (Play! Scala, ElasticSearch, AngularJS)
  \item{} Maintenir le haut niveau de qualité et de tests présents
  \item{} Participer aux ateliers d'architecture et de cadrages techniques
  \item{} Echanger autour de bonnes pratiques avec les équipes de développement
\end{itemize}
\begin{figure}[h]
  \begin{center}
    \hspace*{-1in}
    \includegraphics[width=1\textwidth]{Pictures/EspaceRecruteur.png}
  \end{center}
\end{figure}

%-------------------------------------------------------------------------------
\subsection{Méthode de management}
\label{sub:Méthode de management}
%-------------------------------------------------------------------------------
\paragraph{}
L'équipe Cadremploi, sur le projet Espace Recruteur notamment, suit une méthodologie Agile.
Les méthodes de type agile se veulent plus pragmatiques que les méthodes traditionnelles.
Elles impliquent au maximum le client et permettent une grande réactivité à ses demandes.
Les méthodes agiles prônent 4 valeurs fondamentales: l'équipe, qui doit être soudée et dans laquelle la communication est fondamentale, le fonctionnement de l'application, la collaboration avec les clients au delà de la négociation contractuelle ainsi que la flexibilité de la structure du logiciel, acceptant ainsi aisément le changement.
Pour l'équipe Cadremploi, cela se caractérisait par des meetings quotidien devant un tableau rassemblant des tâches demandées par le client écrites sur des post-its, une mise en production des nouveautés développées chaque jeudi, la mise en place de tests unitaires et end-to-end ainsi que des revues de code pour chaque tâche effectuée.
%-------------------------------------------------------------------------------
\paragraph{Le stand-up meeting}
Comme dans de nombreuses méthodes Agiles, la communication au sein de l'équipe passe notamment par un stand-up-meeting quotidien ainsi qu'un affichage mural composé de posts-its représentant les tâches relatives à l'évolution de l'application.
Le système de tickets post-it, des stand-up-meeting était très bien suivi bien que leur pratique ne soit pas un exercice figé.
Le contenu des stand-up-meeting est par exemple passé au cours de mon stage de "Qu'est-ce que j'ai fait?, Qu'est-ce que je vais faire?, Ai-je des problèmes?" à des questions plus pratiques: "Qu'est-ce que j'ai fait de réellement impactant pour le reste de l'équipe? Quels dysfonctionnements ai-je remarqué?".
Finalement, plutôt que l'évolution de chacun dans sa tâche, c'est l'évolution de l'application elle-même qui était alors discutée durant les stands-up, réduisant ainsi le œtemps de l'exercice, ce qui était une nécessité dans cette équipe de 10 personnes.
De plus, suite à un amoncellement de tâches trop important qui produisait un confus global lors de la lecture du tableau, il a été mis en place des règles permettant de limiter le nombre de tickets en cours de traitement, permettant une meilleure organisation de l'équipe.
%-------------------------------------------------------------------------------
\paragraph{Cycle de vie d'une tâche}
Les post-its suivaient un cycle de vie dont les étapes rentraient dans les catégories suivantes:
\begin{itemize}
  \item "Nouveau", il s'agit du statut des tickets qui viennent d'arriver sur le tableau et qui seront donc à traiter prochainement
  \item "En cours", il s'agit d'un ticket qu'un membre de l'équipe est en train de traiter.
  \item "Recette" équipe, lorsqu'une tâche est terminée, elle doit être validée par un membre de l'équipe.
  Ce système de team review permet d'éviter de nombreuses erreurs au cours du développement de l'application: une vérification du code ainsi que des tests écrits sont effectués.
  \item "Recette" PO lorsque le ticket passe l'étape de recette équipe, il arrive en recette PO.
  Le Project Owner s'occupe donc de vérifier que les modifications apportées correspondent bien aux attentes et qu'elles ne perturbent pas les anciennes fonctionnalités.
  Si un problème est rencontré, le PO rédige ses retours et la personne responsable du ticket se charge de corriger les problèmes.
  Le ticket reste dans le même état mais un jeu d'étiquettes "Retours Recette" et "Corrigé" est déposé sur le ticket pour qu'il soit possible de suivre l'évolution des corrections.
  \item "Validé", le PO valide le ticket et les modifications seront mises en production sous peu
  \item "En attente", le ticket ne peut être traité pour le moment; une synchronisation avec une autre équipe est nécessaire ou bien des questions restent sans réponse.
\end{itemize}
Ainsi les membres de l'équipe choisissent arbitrairement un ticket qu'il traitent en collant une étiquette, un autocollant les représentant dessus et déplacent le ticket dans la catégorie à laquelle il correspond.
Les tickets validés sont généralement mis en production chaque semaine, le jeudi, à l'exception de quelques mises en productions (MEP) exceptionnelles.
\paragraph{Organisation du tableau}
Le tableau sur lequel sont affichés les tickets est plutôt grand et plusieurs méthodes sont utilisées pour permettre à l'équipe de s'y retrouver rapidement.
En effet, en plus des colonnes de catégories dans lesquelles sont rangés les tickets, un code couleur permettait de différencier les tâches techniques des tâches fonctionnelles, ainsi que les tâches traitant de parties de Cadremploi autre que l'espace recruteur (projet de refonte de la Home, ...).
De plus, de petites étiquettes, en plus de celles représentant les membres de l'équipe, sont placées sur les tickets et permettent d'indiquer facilement si un des posts-its représente une tâche qui doit être mise en production dans la semaine (Cible MEP), ou si des retours suite à une recette (équipe ou PO) sont à traiter.
\paragraph{Référencement des tâches}
Enfin, les post-its représentant nos tâches rédigées sur un tableau se trouvent aussi référencées sur l'outil Redmine, une application web de gestion de projets, qui stocke toutes les informations concernant une tâche donnée.
Ce stockage double de l'information est un peu lourd; en effet, il est nécessaire de déplacer le ticket sur le tableau et d'effectuer le même changement en ligne sur Redmine.
Le point positif est qu'il permet à toute l'équipe, même en télétravail, d'accéder facilement à l'information concernant une tâche, notamment aux différentes discussions ayant eues lieu concernant cette tâche.
Il est cependant nécessaire de conserver l'affichage au tableau puisqu'il est bien plus visuel et permet à l'équipe d'échanger bien plus facilement que si le support était uniquement digital.
